\section{Erd\H{o}s Problem \#199}
\subsection*{1. FORMAL RESTATEMENT}
Suppose $A\subset\mathbb R$ contains no nonconstant three-term arithmetic progression. Must $\mathbb R\setminus A$ contain a set $\{x+nd:n\ge0\}$ with $d\ne0$?
\subsection*{2. KNOWN STATUS AND ATTEMPT SCOPE}
The current database attributes a negative answer to Baumgartner's 1975 Hamel-basis construction. The original publisher record was checked, but its full proof was not accessible in this attempt. We do not present the theorem's statement as our own proof. Instead we prove two substantive boundary results: the assertion is true for every Lebesgue measurable $A$, whereas the analogous assertion fails over $\mathbb Q$. The missing step is the unrestricted real counterexample.
\subsection*{3. MEASURABLE PROGRESSION-FREE SETS HAVE MEASURE ZERO}
Suppose a measurable progression-free $A$ had positive measure. Some bounded measurable $B\subseteq A$ would have $0<\mu(B)<\infty$. Translation is continuous in $L^1$, so for all sufficiently small nonzero $t$,
\[
 \mu(B\mathbin\triangle(B-t))<\frac{\mu(B)}4,
 \qquad
 \mu(B\mathbin\triangle(B-2t))<\frac{\mu(B)}4.
\]
Therefore
\[
 \mu\bigl(B\cap(B-t)\cap(B-2t)\bigr)
 \ge\mu(B)-\mu(B\setminus(B-t))-\mu(B\setminus(B-2t))>0.
\]
A point in this intersection gives $x,x+t,x+2t\in B$, a contradiction. The translation-continuity fact follows, for example, by approximating the indicator of a finite-measure measurable set in $L^1$ by a finite linear combination of interval indicators and checking interval translations directly.

Now fix any nonzero real $d$. Since $A$ is null,
\[
 E_d=\bigcup_{n\ge0}(A-nd)
\]
is null. Every $x\notin E_d$ has $x+nd\notin A$ for all $n\ge0$. Thus for each prescribed nonzero $d$, almost every starting point gives an infinite progression in the complement. In particular any counterexample to the original question must be nonmeasurable. This rules out every countable or Borel candidate as well.
\subsection*{4. A COUNTEREXAMPLE OVER THE RATIONALS}
There are only countably many progressions $P=\{x+nd:n\ge0\}$ with rational $x$ and nonzero rational $d$. Enumerate them, with repetitions if desired. Inductively maintain a finite progression-free set $F$. In choosing a new point $y$, the only ways to create a progression with two old points $a,b\in F$ are
\[
 y=(a+b)/2,\qquad y=2a-b,\qquad y=2b-a.
\]
Together with $F$, these form a finite forbidden set. The next progression is infinite, so choose a point on it outside that finite set. The union $A_{\mathbb Q}$ is progression-free, because any prohibited triple would appear at a finite stage, and it meets every rational progression. Hence $\mathbb Q\setminus A_{\mathbb Q}$ contains no infinite rational progression.

This does not give a counterexample in $\mathbb R$: $A_{\mathbb Q}$ is countable and the preceding measure argument applies to it. Nor can one replace countable recursion by an uncountable recursion without further proof: a countable progression may then be exhausted by the previously forbidden points.
\subsection*{5. FINAL}
\textbf{UNRESOLVED in this reconstruction, although the original problem is known to be disproved.}
We have established the entire measurable case, its null-set strengthening for every fixed step, and the contrasting rational construction. Reconstructing Baumgartner's coherent construction across all real rational-affine lines remains necessary for a full independent resolution. The estimate below assesses this attempt, not the literature's settled status.
\subsection*{6. SOURCES}
Current statement and historical attribution: \texttt{https://www.erdosproblems.com/199}. J. E. Baumgartner, \emph{Partitioning vector spaces}, J. Combinatorial Theory Ser. A 18 (1975), 231--233, \texttt{https://doi.org/10.1016/0097-3165(75)90016-3}. This attempt uses elementary measure theory and finite avoidance, not an unverified extraction of that paper's proof.
\section{COMPLETION ESTIMATE}
\textbf{COMPLETION: 40\%}
